
\section{System Architecture}
\label{sec:system-architecture}

MLguide is structured as a three-tier web application consisting of a 
frontend, a backend, and a database layer, deployed as Docker containers 
\parencite{docker}. Docker was chosen because containerisation makes the 
system easy to run across different environments without managing 
dependencies manually, which is important for a research artefact that 
needs to be shared and tested by others.

Figure~\ref{fig:system-diagram-components} shows the system architecture of MLguide,
with the main components as different
services, and how they communicate. 
The components are grouped
into application services, data services, and external services. The
application and data services run in Docker containers, while the
external services are accessed over the network. The application
services consist of the frontend and the backend. 
The data services consist of the ontology-based knowledge graph and the SQL database.

\begin{figure}[H]
    \centering
    \includegraphics[width=\textwidth]{Figures/06-implementation/system-diagram-components.pdf}
    \caption[MLguide system architecture]{Overview of the MLguide system architecture, showing the application, data, 
    and external services and how they communicate.}
    \label{fig:system-diagram-components}
\end{figure}

The frontend is implemented in Streamlit 
\parencite{streamlit}, which was chosen because it allows interactive 
web applications to be built directly in Python, without writing HTML or 
JavaScript. This made it well suited for rapid development of a 
data-driven prototype. The frontend communicates with the backend through 
an Application Programming Interface (API) over the Hypertext Transfer Protocol (HTTP). The backend is implemented in FastAPI \parencite{fastapi}, 
which provides automatic API documentation and built-in data validation,
and contains the core application logic.

The backend connects to two data stores. GraphDB \parencite{graphdb} stores 
the ontology-based knowledge graph, and is queried using SPARQL. 
It was chosen because it is designed specifically for RDF data and 
supports SPARQL, making it a natural fit for an RDF-based knowledge graph. 
PostgreSQL \parencite{postgresql} stores article abstracts 
and vector embeddings using the pgvector extension \parencite{pgvector}. 
This allows embeddings and relational data to be stored and queried in the 
same database, avoiding the need for a separate vector database. 
The system also integrates with two external services: the GitHub API and 
Google Colab, which are used to upload and open generated notebooks.

MLguide follows a layered architecture, in which each layer only depends 
on the layer immediately beneath it 
\parencite[157--158]{sommerville2011software_engineering}. This keeps concerns 
separated across the system, so that individual layers can be modified or 
replaced without affecting the rest of the application. The layered 
approach also supports incremental development, because each layer can be 
made available as it is completed without waiting for the full system to 
be built \parencite[157--158]{sommerville2011software_engineering}. This is 
consistent with the development process described in 
Section~\ref{sec:incremental-and-iterative-development}.

Figure~\ref{fig:mlguide-layered-architecture} shows the layered architecture
of MLguide including the internal layer
structure of the frontend and backend. Both are organised as a
stack of layers with clear responsibilities.
The frontend and backend also have shared definitions used across several layers. In
the frontend, the data models define the request and response shapes
exchanged with the backend. The templates hold the starter notebook 
templates used in notebook generation. 
In the backend, the schemas define the
request and response formats used across the layers.

\begin{figure}[H]
    \centering
    \includegraphics[width=\textwidth]{Figures/06-implementation/mlguide-layered-architecture.pdf}
    \caption[MLguide layered architecture]{Layered architecture of MLguide.}
    \label{fig:mlguide-layered-architecture}
\end{figure}

In the backend, the router layer 
handles incoming API requests and passes them to the service layer, which 
contains the core application logic. The persistence layer is responsible for 
all database access and isolates data retrieval from the rest of the 
application. In the frontend, the presentation layer handles rendering 
and user interaction. The application logic layer coordinates calls to 
the backend and external services. The application logic layer delegates most logic
to the backend, so the majority of 
modules in this layer are only thin wrappers around the integration layer.  
The clearest exceptions are the notebook template services, which contain local 
logic for selecting and rendering starter notebooks.
The integration layer holds the clients that handle the actual
communication with the backend and external services.

Figure~\ref{fig:package-diagram-full} shows the package diagram for
MLguide, including the main modules within each layer of the frontend
and backend.

\begin{figure}[H]
    \centering
    \includegraphics[width=\textwidth]{Figures/06-implementation/package-diagram-full.pdf}
    \caption[MLguide package diagram]{Package diagram of MLguide.}
    \label{fig:package-diagram-full}
\end{figure}

The package diagram shows that the recommendation process is the core 
of the system, with multiple services in the backend application logic layer.
Alongside the recommendation flow, the system
includes two supporting features that have not been introduced yet:
user services and feedback handling. User services is a lightweight
feature covering authentication and saved searches, allowing users to
store and reload previous searches with problem descriptions. Feedback
handling allows users to rate recommended methods. The ratings are stored
in PostgreSQL and can be used to adjust the recommendation process over time.
The feedback logic is further described as part of the 
recommendation process in Section~\ref{sec:recommendation-process}. 
Both user and feedback concerns follow the same layered structure as the recommendation logic, 
with separate modules in the router, service, and persistence layers of 
the backend, and corresponding wrappers in the frontend application logic 
layer.
The modules marked as wrappers in the frontend contain little logic of
their own, as described above.
